\bookchapter{제58장: 빠른 참조 매뉴얼과 리소스}

\chapterintro{이번 장은 독립적으로 쓸 수 있는 빠른 참조 매뉴얼이에요. 어떤 문법이 가물가물하거나 빠르게 결정해야 할 때 바로 여기로 넘어오세요. “주머니 카드”처럼 쓰길 권해요. 자주 쓰는 코드 요약, 자주 묻는 질문, 학습 리소스까지 한 페이지도 빠짐없이요.}

\booksection{자주 쓰는 코드 빠른 참조 카드}

\begin{center}
\begin{tabularx}{\linewidth}{|X|X|}
\hline
\textbf{무엇을 하고 싶나} & \textbf{핵심 코드} \tabularnewline
\hline
변경 가능한 상태 하나 기억하기 & var x by remember \{ mutableStateOf(초깃값) \} \tabularnewline
화면을 돌려도 상태 유지 & rememberSaveable \{ mutableStateOf(초깃값) \} \tabularnewline
세로/가로 배치 & \texttt{Column \{ \}}/\texttt{Row \{ \}} \tabularnewline
지연 로딩 리스트 & LazyColumn \{ items(리스트) \{ ... \} \} \tabularnewline
화면 이동 & navController.navigate("라우트") \tabularnewline
코루틴에서 기다리는 작업 처리 & \texttt{viewModelScope.\allowbreak launch \{ .\allowbreak .\allowbreak .\allowbreak  \}} \tabularnewline
데이터 스트림 관찰 & \texttt{val list by flow.\allowbreak collectAsState()} \tabularnewline
작은 설정 저장 & context.dataStore.edit \{ it[KEY] = 값 \} \tabularnewline
구조화된 데이터 저장 & Room 3종 세트: @Entity / @Dao / @Database \tabularnewline
네트워크로 데이터 가져오기 & Retrofit 인터페이스 + suspend 함수 \tabularnewline
알림 보내기 & NotificationChannel + NotificationCompat.Builder \tabularnewline
로그 남기기 & \texttt{Log.\allowbreak d("TAG", "message")} \tabularnewline
\hline
\end{tabularx}
\end{center}

\begin{center}\small\color{gray}안드로이드 개발 빠른 참조 카드(2026년 8월판)\end{center}

\booksection{자주 묻는 질문 FAQ}

Q1: 프로그래밍을 전혀 모르는데, 안드로이드 개발을 배우려면 얼마나 걸리나요?

이 책의 순서를 따라가 보세요. 1-2주 차에 환경을 설치하고 Kotlin 기초를 배우고, 3-4주 차에 Compose를 배워 화면을 만들고, 5-6주 차에 가계부 프로젝트를 완성하고, 7-8주 차에 날씨 프로젝트를 끝내요. 두 달이면 제로에서 작은 앱을 혼자 만들 수 있는 수준까지 갑니다. 조건은 매일 1-2시간 투자하고, 책만 보지 말고 직접 코드를 타이핑하는 것이에요.

Q2: Java는 따로 배워야 하나요?

아니요. 새 프로젝트, 새 튜토리얼, 새 채용 공고는 대부분 Kotlin 기준이에요. “십 년 된 오래된 프로젝트를 유지보수”할 때에만 많은 Java를 만나요. Kotlin을 익힌 뒤 업무에 필요하면 Java를 보충하는 건 빨라요. 둘 다 같은 가상 머신 생태계 위에 있으니까요.

Q3: 에뮬레이터일까요 실제 기기일까요?

개발과 디버깅은 에뮬레이터가 편해요(스냅샷, 여러 기종). 하지만 검수는 반드시 실제 기기에서 해야 해요. 촉감, 성능, 실제 네트워크, 센서는 에뮬레이터가 대체할 수 없으니까요. 여건이 된다면 중급형 실제 기기 한 대를 테스트 기기로 사는 것이 가장 가성비 좋은 투자예요.

Q4: 안드로이드 개발은 취업이 잘 되나요?

2026년 시세: 단순히 “화면을 만들 줄 아는” 초급 포지션은 경쟁이 치열하지만, “Kotlin을 알고 + 아키텍처를 다루고 + 출시한 작품이 있는” 개발자는 여전히 인기가 많아요. 배우는 시기의 여러분에게 최고의 이력서는 스토어에 올려 둔 자신이 만든 앱이에요.

Q5: 오류에 막혔다면 어떻게 하나요?

황금 3단계: (1) 오류 첫 줄을 그대로 StackOverflow/검색 엔진에 검색; (2) 공식 문서의 해당 페이지 확인; (3) 그래도 안 되면 “기대한 동작 + 실제 현상 + 오류 전문 + 관련 코드”를 묶어 AI나 커뮤니티에 질문. 대부분의 문제는 첫 단계에서 해결돼요.

\booksection{추천 학습 리소스}

\begin{center}
\begin{tabularx}{\linewidth}{|X|X|X|}
\hline
\textbf{종류} & \textbf{리소스} & \textbf{추천 이유} \tabularnewline
\hline
공식 & developer.android.com & 안드로이드 공식 문서, Compose/Kotlin 튜토리얼이 모두 들어 있어요. 제일 권위 있는 곳 \tabularnewline
공식 & Kotlin 공식 사이트(kotlinlang.org) & 언어 공식 문서 + 온라인 연습장 \tabularnewline
공식 & Google Codelabs & 손부터 잡아 주는 단계별 실습 과정, 초보자 친화도 만점 \tabularnewline
커뮤니티 & Stack Overflow(android 태그) & 오류로 곤란할 때 가장 먼저 찾는 곳 \tabularnewline
커뮤니티 & 쥐진(Juejin) / CSDN 안드로이드 코너 & 중국어 실전 글이 많아요. 새것과 오래된 것을 가려 보세요(먼저 날짜 확인) \tabularnewline
동영상 & Google Developers 공식 채널 & 매년 열리는 개발자 컨퍼런스 신기술 강연 \tabularnewline
함께 보기 & 제7편 〈AI 코딩 실전편〉(37\textasciitilde{}47장) & AI로 안드로이드 개발을 보조하는 완전 노하우 \tabularnewline
\hline
\end{tabularx}
\end{center}

\begin{center}\small\color{gray}안드로이드 개발 학습 리소스 목록(2026년 8월)\end{center}

\booksection{Gradle 자주 쓰는 명령어 빠른 참조}

평소에는 버튼 클릭으로 충분하지만, 명령줄이 때로는 더 효율적이에요(특히 AI 도구나 CI 파이프라인과 함께 쓸 때):

\begin{center}
\begin{tabularx}{\linewidth}{|X|X|}
\hline
\textbf{명령어} & \textbf{역할} \tabularnewline
\hline
\texttt{.\allowbreak /\allowbreak gradlew assembleDebug} & 디버그용 APK 빌드(Windows에서는 ./ 제거) \tabularnewline
\texttt{.\allowbreak /\allowbreak gradlew assembleRelease} & 릴리스 빌드(서명 설정 필요) \tabularnewline
\texttt{.\allowbreak /\allowbreak gradlew test} & 단위 테스트 전부 실행 \tabularnewline
\texttt{.\allowbreak /\allowbreak gradlew lint} & 정적 검사: 잠재적 버그와 관례에 어긋나는 코드 찾기 \tabularnewline
\texttt{.\allowbreak /\allowbreak gradlew clean} & 빌드 산출물 비우기("재시작 만능법"의 컴파일 버전) \tabularnewline
\texttt{.\allowbreak /\allowbreak gradlew :app: dependencies} & 의존성 트리 확인(버전 충돌 추적 필살 도구) \tabularnewline
\hline
\end{tabularx}
\end{center}

\begin{center}\small\color{gray}Gradle 자주 쓰는 명령어 빠른 참조\end{center}

\begin{tcbinfo}{팁: 빌드 문제가 생기면 먼저 clean}
“어제는 됐는데 오늘은 안 돼” 같은 기이한 현상은 열에 여섯일곱이 빌드 캐시 꼬임 탓이에요. ./gradlew clean 후에 다시 빌드하면 미스터리 문제의 태반이 해결돼요. 그래도 안 되면 Android Studio를 재시작해 보세요. 이때 File Invalidate Caches와 함께 쓰면 효과가 배가돼요.

\end{tcbinfo}

\booksection{맺음말: 첫 앱에서 N번째 앱으로}

\booksection{맺음말: 첫 앱에서 N번째 앱으로}

돌아보면, 여러분은 “Activity가 뭔지도 몰랐던” 상태에서 환경을 설치하고, Kotlin을 배우고, Compose로 화면을 그리고, 앱에 데이터베이스와 네트워크를 연결한 뒤, 테스트·패키징·스토어 등록까지 해냈어요. 이 길에서 여러분은 이미 0에서 1로 가는 가장 어려운 구간을 걸어 넘은 거예요.

기술의 세계는 변화가 빨라요. 프레임워크는 업그레이드되고, 도구는 세대를 거듭하고, AI는 점점 더 많은 “키보드 두드리는” 일을 대신 하죠. 하지만 한 가지는 낡지 않아요. 막연한 아이디어를 명확한 단계로 쪼개고 한 걸음씩 실현해 나가는 능력이에요. 도구는 증폭기이고, 사고방식은 기초예요. 증폭기는 바꿀 수 있지만 기초는 흔들리면 안 돼요.

더 이상 "휴대폰 잘 쓰는 사람"이 아니에요. 머릿속 아이디어를 3억 대의 안드로이드 기기에서 실제로 돌아가는 앱으로 만들 수 있어요. 오늘부터 키보드로 만들어 보세요 — 이것이 안드로이드 개발자의 낭만이에요.

\booksection{참고 문헌}

\begin{enumerate}
  \item Android Developers. (2026). Android Developer Guides and Jetpack Compose Documentation. developer.android.com.
  \item Kotlin Foundation. (2026). Kotlin Language Documentation. kotlinlang.org.
  \item Android Open Source Project. (2026). Android Platform Version History and API Levels.
  \item Google Play. (2026). Play Console Help: Publishing and App Signing.
  \item AndroidX Release Notes. (2026). Room / DataStore / Navigation / WorkManager releases.
  \item Square. (2026). Retrofit Documentation. github.com/square/retrofit.
\end{enumerate}
